\section{Eksperyment kontrfaktyczny BDP}

Scenariusz BDP jest kontrfaktycznym rozszerzeniem istniejącego punktu bazowego, a nie nową kalibracją modelu. W eksperymencie zmieniane są dwa elementy: miesięczny transfer do gospodarstw domowych oraz siła, z jaką transfer podnosi płacę rezerwową. Oba kanały są wyłączone w wariancie bez BDP, dzięki czemu scenariusze BDP można interpretować jako odchylenia od tej samej ścieżki startowej.

Transfer jest traktowany jako wydatek społeczny sektora publicznego. Dlatego oddziałuje jednocześnie na dochód gospodarstw domowych, deficyt, dług, popyt, inflację, koszt kapitału i bilanse firm. To jest kluczowe dla hipotezy: BDP może zwiększać presję automatyzacyjną, ale ta presja musi zostać sfinansowana przez firmy w środowisku zmieniającego się kosztu kapitału.

Formuła transferu gospodarstwa domowego ma postać:
\begin{equation}
  Transfer_h = Children_h \cdot 800 + BDP.
\end{equation}
Kanał płacy rezerwowej zapisano jako:
\begin{equation}
  ReservationWage = MinWage + \lambda \cdot BDP.
\end{equation}
W wariancie centralnym przyjęto \(\lambda = 0{,}5\). Jest to założenie behawioralne: gospodarstwo domowe otrzymuje pełny BDP, ale tylko część transferu podnosi jego płacę rezerwową. Dla BDP równego zero model pozostaje wariantem bazowym bez dodatkowego transferu. W analizie odporności siła kanału płacy rezerwowej jest zmieniana dla \(\lambda = 0\), \(0{,}25\), \(0{,}5\), \(0{,}75\) i \(1\).

Eksperyment objął poziomy \(0\), \(500\), \(1000\), \(1500\), \(2000\), \(2500\) i \(3000\) PLN miesięcznie na agenta reprezentującego gospodarstwo domowe. Horyzont wynosił 60 miesięcy, a każdy poziom uruchomiono dla 10 ziaren losowania. Szczegółowa ścieżka odtworzenia, commit bazowy, plik zmian i komendy znajdują się w aneksie replikacyjnym.

Jednostka transferu wymaga doprecyzowania. Poziom BDP 2000 PLN nie oznacza literalnie 2000 PLN miesięcznie na każdego mieszkańca Polski. Oznacza miesięczny transfer na agenta reprezentującego gospodarstwo domowe w skalowanej populacji modelu. Dlatego przeliczenie na język decydenta publicznego powinno wykorzystywać dodatkowy miesięczny strumień transferów społecznych względem wariantu bez BDP, podzielony przez populację referencyjną 38 mln osób, a nie proste dzielenie liczby mieszkańców przez liczbę agentów.

Pełny zbiór wyników modelu jest szerszy niż zakres tej pracy. W analizie wykorzystuję tylko wskaźniki odpowiadające kanałom hipotezy: skalę transferu, deficyt i dług, koszt kapitału, kurs i rachunek bieżący, inwestycje prywatne, adopcję AI/hybrid, demografię firm oraz kontrolnie wybrane miary gospodarstw domowych i sektora bankowego.

\begin{table}[htbp]
\centering
\scriptsize
\caption{Mapa kanałów transmisji scenariusza BDP}
\label{tab:transmission-channels}
\begin{tabularx}{\textwidth}{@{}>{\raggedright\arraybackslash}p{0.16\textwidth}>{\raggedright\arraybackslash}p{0.32\textwidth}>{\raggedright\arraybackslash}p{0.18\textwidth}>{\raggedright\arraybackslash}X@{}}
\toprule
\textbf{Kanał} & \textbf{Mechanizm ekonomiczny} & \textbf{Wskaźniki} & \textbf{Interpretacja wyniku w wariancie centralnym} \\
\midrule
Płaca rezerwowa &
Transfer zwiększa dochód niezależny od pracy; część transferu podnosi płacę rezerwową przez parametr \(\lambda\). &
płaca rynkowa, płaca zatrudnionych, bezrobocie &
Kanał jest aktywny mechanicznie, ale płaca rynkowa reaguje silniej dopiero przy wysokim BDP; sam kanał nie tłumaczy lokalnego maksimum adopcji przy \texttt{--bdp 2000}. \\
\addlinespace
Popyt i dochód gospodarstw &
BDP zwiększa dochód gospodarstw domowych, lecz część impulsu może trafić do oszczędności, importu i bilansów zamiast do krajowego popytu finalnego. &
dochód HH, mediana konsumpcji, PKB, transfery społeczne &
Dochód gospodarstw rośnie, ale mediana konsumpcji i PKB reagują słabiej; kanał popytowy wspiera firmy tylko częściowo. \\
\addlinespace
Fiskus i dług &
Transfer zwiększa wydatki publiczne, deficyt oraz dług publiczny względem PKB. &
deficyt/PKB, dług/PKB, rentowność długu, obsługa długu &
To najsilniejszy kanał monotoniczny: przy wysokim BDP presja fiskalno-dłużna zaczyna dominować nad pozytywnym impulsem popytowym. \\
\addlinespace
Koszt kapitału &
Inflacja, reakcja NBP, dług publiczny i premia rynkowa wpływają na koszt finansowania inwestycji firm. &
inflacja, stopa NBP, obligacje, dług korporacyjny, inwestycje/PKB &
Stopa NBP rośnie umiarkowanie, ale rentowności obligacji i koszt długu korporacyjnego rosną wyraźniej; inwestycje prywatne spadają. \\
\addlinespace
Sektor zewnętrzny &
Wyższy popyt i importochłonność inwestycji wpływają na import, kurs PLN/EUR i rachunek bieżący. &
kurs, import, rachunek bieżący, ceny importu &
PLN słabnie, import rośnie, a rachunek bieżący pogarsza się; część impulsu popytowego odpływa przez kanał zewnętrzny. \\
\addlinespace
Kredyt i banki &
Finansowanie inwestycji zależy od jakości portfela kredytowego, płynności i adekwatności kapitałowej banków. &
NPL, CAR, LCR, NSFR, upadłości banków &
Nie występuje kryzys bankowy; ograniczenie automatyzacji przy wysokim BDP wynika raczej z kosztu i opłacalności finansowania niż z upadłości banków. \\
\addlinespace
Automatyzacja i firmy &
Firmy porównują koszt pracy, popyt, koszt kapitału i dostępność finansowania; automatyzacja wymaga sfinansowania CapEx. &
adopcja AI/hybrid, inwestycje, nowe firmy i firmy kończące działalność &
Adopcja AI/hybrid rośnie do \texttt{--bdp 2000}, po czym spada; więcej transferu nie oznacza większej automatyzacji, jeśli spada zdolność finansowania inwestycji. \\
\bottomrule
\end{tabularx}
\end{table}


Wyniki prezentuję jako tabele syntetyczne i wykresy kanałów transmisji. Pełne nazwy kolumn, pliki wynikowe i procedura generowania artefaktów pozostają częścią aneksu replikacyjnego.
